\documentclass[a4paper]{article}

% Basic packages
% Español (México) con XeLaTeX: fontspec para fuentes Unicode, polyglossia
% para la división silábica y los nombres del documento en la variante mexicana.
\usepackage{fontspec}
\usepackage{polyglossia}
\setdefaultlanguage[variant=mexican]{spanish}
% Los nombres chinos van como «pinyin (original)»: xeCJK da a los caracteres
% Han su propia fuente; sin ella XeLaTeX los omite con un aviso.
% FandolSong no cubre algunos caracteres de nombres (U+73FA); WenQuanYi Zen Hei sí.
\usepackage[AutoFallBack=true]{xeCJK}
\setCJKmainfont{FandolSong-Regular.otf}
\setCJKfallbackfamilyfont{\CJKrmdefault}{WenQuanYi Zen Hei}
% polyglossia no traduce los nombres de listings.
\AtBeginDocument{\providecommand{\lstlistingname}{}\renewcommand{\lstlistingname}{Listado}%
  \providecommand{\lstlistlistingname}{}\renewcommand{\lstlistlistingname}{Índice de listados}}
\usepackage{amsmath, amssymb}
\usepackage{graphicx}
\usepackage[margin=2.5cm]{geometry}
\usepackage[most]{tcolorbox}
\usepackage{etoolbox}
\usepackage{listings}
\usepackage{hyperref}
\usepackage{booktabs}
\usepackage{subcaption}
\usepackage{float}

% Highlight boxes
\newtcolorbox{knowledgebox}[1]{
    enhanced, colback=blue!5!white, colframe=blue!75!black, colbacktitle=blue!75!black,
    coltitle=white, fonttitle=\bfseries, title={#1}, attach boxed title to top left={yshift=-2mm, xshift=2mm},
    boxrule=1pt, sharp corners
}

\newtcolorbox{importantbox}[1]{
    enhanced, colback=yellow!10!white, colframe=yellow!80!black, colbacktitle=yellow!80!black,
    coltitle=black, fonttitle=\bfseries, title={#1}, sharp corners
}

\newtcolorbox{warningbox}[1]{
    enhanced, colback=red!5!white, colframe=red!75!black, colbacktitle=red!75!black,
    coltitle=white, fonttitle=\bfseries, title={#1}, sharp corners
}

% Listings
\lstset{
    language=Python,
    basicstyle=\ttfamily\small,
    keywordstyle=\color{blue},
    stringstyle=\color{red!60!black},
    commentstyle=\color{green!60!black},
    breaklines=true,
    frame=single,
    numbers=left,
    numberstyle=\tiny\color{gray},
    captionpos=b,
    extendedchars=true
}

% Front-page metadata
\newcommand{\notetitle}{How We Scaled Kimi K2.5}
\newcommand{\noteauthors}{Elaboradas a partir de materiales públicos del curso}
\newcommand{\notedate}{2026}
\newcommand{\videochannel}{NVIDIA GTC 2026}
\newcommand{\videopublishdate}{2026}
\newcommand{\videoduration}{39 min}
\newcommand{\videourl}{}
\newcommand{\videocoverpath}{cover.jpg}
\newcommand{\repourl}{https://github.com/hqhq1025/ai-course-notes}


% Footer with repo info
\usepackage{fancyhdr}
\pagestyle{fancy}
\fancyhf{}
\fancyfoot[C]{\thepage}
\fancyfoot[R]{\footnotesize\color{gray}\href{\repourl}{AI Course Notes}}
\renewcommand{\headrulewidth}{0pt}
\renewcommand{\footrulewidth}{0pt}

\begin{document}

\begin{titlepage}
\centering
{\Large Notas de la conferencia\par}
\vspace{1.2cm}
{\huge\bfseries \notetitle\par}
\vspace{0.5cm}
{\large Yang Zhilin (杨植麟) (Moonshot AI / Kimi) --- GTC 2026 Keynote\par}
\vspace{0.8cm}
{\large \noteauthors\par}
\vspace{0.3cm}
{\large \notedate\par}
\vspace{1.2cm}

\ifdefempty{\videocoverpath}{}{
\includegraphics[width=0.82\textwidth,height=0.45\textheight,keepaspectratio]{\videocoverpath}\par
}

\vfill
\begin{tcolorbox}[width=0.9\textwidth, colback=black!2!white, colframe=black!60, sharp corners]
\textbf{Ponente}: Yang Zhilin (杨植麟) (Moonshot AI / Kimi)\par
\textbf{Ocasión}: NVIDIA GTC 2026\par
\textbf{Duración del video}: \videoduration\par
\textbf{Enlace del video}: 
\ifdefempty{\videourl}{}{\textbf{Enlace del video}: \href{\videourl}{\nolinkurl{\videourl}}\par}\par
\textbf{Página del proyecto}: \href{\repourl}{\nolinkurl{\repourl}}\par
\end{tcolorbox}
\end{titlepage}

\tableofcontents
\newpage

%% --- Inicio del contenido --- %%

\section{Introducción: la visión de los modelos abiertos}

Yang Zhilin (杨植麟) expuso, en su conferencia en el GTC 2026, los avances más recientes de Moonshot AI (Kimi) en el campo de los modelos abiertos. Citó un slide que Jensen Huang mostró en el CES, en el que se señala que los modelos abiertos están cerrando rápidamente la brecha con los modelos de código cerrado y están alcanzando el nivel de vanguardia.

\begin{importantbox}{Idea central}
Los modelos abiertos no solo deben ser «abiertos», sino también «excelentes» (Open models cannot be just open --- they have to be great). El objetivo del equipo de Kimi es que la inteligencia llegue, a través de modelos de código abierto, a cada rincón del mundo, y al mismo tiempo competir en desempeño con los modelos de código cerrado.
\end{importantbox}

La conferencia se desarrolla en torno a tres dimensiones de scaling:
\begin{enumerate}
    \item \textbf{Token Efficiency}: mediante mejores arquitecturas y optimizadores, se mejora la eficiencia de aprendizaje de cada token
    \item \textbf{Long Context}: se amplía la longitud del contexto para dar soporte a tareas de agent más complejas
    \item \textbf{Agent Swarms}: múltiples agents colaboran en paralelo para completar tareas complejas
\end{enumerate}

\begin{knowledgebox}{De scaling a agent: una perspectiva unificada}
Yang Zhilin expresa estas tres dimensiones de manera unificada en el lenguaje de agent: token efficiency corresponde a un conocimiento previo más fuerte (la búsqueda del agent RL es más eficiente); long context corresponde a un tiempo de ejecución del agent más largo (puede ejecutarse durante días o incluso semanas); agent swarms es, por su parte, la nueva dimensión de la paralelización. El objetivo final es contar con un clúster compuesto por agents con un contexto extremadamente largo y un conocimiento previo extremadamente fuerte.
\end{knowledgebox}

\section{Token Efficiency: el optimizador Muon y la estabilidad del entrenamiento}

\subsection{Por qué la token efficiency es tan importante}

Yang Zhilin (杨植麟) parte de la clásica Kaplan scaling law para enfatizar que la token efficiency no solo concierne a la eficiencia, sino también al \textbf{límite superior de la inteligencia}.

\begin{importantbox}{Token Efficiency = incremento del límite superior de la inteligencia}
Supongamos que se cuenta con 50 billones de tokens de alta calidad; si un nuevo optimizador logra duplicar la token efficiency, eso equivale a obtener «de la nada» el efecto de 100 billones de tokens. En una época en la que el volumen de datos de alta calidad se acerca a su límite (data wall), incrementar la token efficiency significa directamente elevar el techo de la inteligencia.
\end{importantbox}

\subsection{El optimizador Muon}

Muon es un \textbf{optimizador de segundo orden} cuya idea central es transformar cada actualización de gradiente de modo que las distintas componentes de los parámetros queden ortogonales entre sí. En comparación con el tradicional AdamW, Muon mejora de manera notable el desempeño con la misma cantidad de parámetros y de tokens de entrenamiento.

El equipo de Kimi fue el primero en demostrar que el optimizador Muon puede escalar al entrenamiento de LLM a gran escala; las técnicas centrales incluyen:
\begin{itemize}
    \item \textbf{Estrategia de decay}: fundamental para escalar a modelos más grandes
    \item \textbf{Consistencia de RMS}: se introduce un coeficiente ajustable para asegurar que la actualización de RMS de Muon sea comparable con la de Adam
    \item \textbf{Implementación distribuida}: se fragmenta el estado del optimizador entre el data parallel group, lo que logra una implementación distribuida eficiente de Muon
\end{itemize}

\subsection{QK-Clip: cómo resolver la inestabilidad del entrenamiento}

Al escalar Muon a un modelo de 1 billón de parámetros, el equipo se topó con un problema de inestabilidad en el entrenamiento: los max logits se disparaban rápidamente por encima de 1000 (el valor normal es de alrededor de 50 a 100), lo que provocaba que el entrenamiento divergiera.

\begin{importantbox}{La técnica QK-Clip}
Para cada attention head, se calcula el max logit durante el forward pass y luego se calcula un factor divisor que se aplica a las proyecciones de query y de key, con lo que se restringe el valor máximo dentro de un rango dado. Los experimentos muestran que:
\begin{itemize}
    \item Las curvas de entrenamiento antes y después de aplicar QK-Clip coinciden de manera estricta, es decir, no afecta la convergencia del entrenamiento
    \item Una vez que el max logit alcanza el umbral (por ejemplo, 100), queda restringido de manera efectiva y después disminuye de forma natural
\end{itemize}
Esto permitió que el equipo de Kimi completara con éxito el primer entrenamiento a gran escala con Muon en la historia (1 billón de parámetros).
\end{importantbox}

\begin{warningbox}{La estabilidad del entrenamiento no se puede ignorar}
El problema de logit explosion, que rara vez aparece en el entrenamiento tradicional con Adam, puede amplificarse al usar un optimizador de segundo orden. Cambiar el optimizador sin más, sin considerar mecanismos de estabilidad, puede provocar que el entrenamiento fracase por completo.
\end{warningbox}

\subsection{Resumen de la sección}

La token efficiency no es solo una optimización de ingeniería, sino también una vía clave para elevar el límite superior de la inteligencia. El optimizador Muon logra duplicar la token efficiency al ortogonalizar las actualizaciones de gradiente; la técnica QK-Clip resuelve el problema de estabilidad en el entrenamiento a gran escala y permite que Muon escale con éxito hasta el billón de parámetros.

\section{Long Context: Kimi Linear y Delta Attention}

\subsection{Por qué el contexto largo es fundamental en la era de los agentes}

Yang Zhilin (杨植麟) citó una gráfica clásica, del nivel de una «joya oculta»: la comparación entre Transformer y LSTM. El Transformer no solo obtiene una loss menor con los mismos parámetros y el mismo número de tokens, sino que, de manera más importante, es capaz de \textbf{mejorar de manera continua su predicción por medio del contexto}: a medida que aumenta el índice de tokens, la loss del Transformer sigue bajando, mientras que la del LSTM se satura después de cierto rango.

\begin{knowledgebox}{La verdadera ventaja del Transformer}
El Transformer no gana solo por ser «más fuerte», sino porque puede aprovechar un contexto más largo para obtener una ganancia de información continua. Esta característica es especialmente crucial en la era de los agentes: un agente necesita ejecutarse durante días o incluso semanas para completar tareas complejas (como escribir un kernel de Linux desde cero), lo cual exige que el modelo procese de manera eficiente una trayectoria extremadamente larga.
\end{knowledgebox}

\subsection{Kimi Delta Attention: decaimiento de grano fino}

La innovación central de la arquitectura Kimi Linear es \textbf{Kimi Delta Attention}, que mejora el mecanismo de memoria recurrente de GDR (Generalized Delta Rule).

La linear attention tradicional usa un factor de decaimiento escalar global (global scalar decay), lo cual obliga a elegir entre «olvidarlo todo» y «recordarlo todo». Kimi Delta Attention introduce un \textbf{factor de decaimiento de grano fino} (fine-grained decay):

$$\alpha \in \mathbb{R}^{d \times d} \quad (\text{matriz diagonal, no un escalar})$$

De este modo, distintos canales pueden tener distintas velocidades de decaimiento:
\begin{itemize}
    \item algunos canales decaen lentamente, con lo que conservan información de largo alcance
    \item algunos canales decaen rápidamente, con lo que olvidan a tiempo y absorben información nueva
\end{itemize}

\subsection{Implementación eficiente: fórmula por chunks e inversión de matrices}

Para paralelizar en las GPU actuales, es necesario convertir la fórmula recurrente en una forma por chunks (chunk-wise). Pero la matriz diagonal $\alpha$ no permite extraer un factor común con la misma facilidad que un escalar, lo cual plantea un enorme desafío de ingeniería.

\begin{importantbox}{Una fórmula paralela exactamente equivalente}
Al introducir una \textbf{operación de inversión de matrices} y un \textbf{factor de decaimiento acumulado}, el equipo reescribió la fórmula recurrente como tres ecuaciones que se pueden calcular en paralelo. No se trata de una aproximación: es una fórmula matemáticamente equivalente en sentido estricto, que logra un paralelismo eficiente sin sacrificar precisión alguna.
\end{importantbox}

\subsection{Comparación de desempeño}

Kimi Linear obtuvo una ventaja integral en comparaciones justas:
\begin{itemize}
    \item \textbf{Tareas de contexto corto} (MMAU): superior a MLA y a GDN
    \item \textbf{Tareas de contexto largo} (Ruler): también superior a las demás variantes, y con mayor eficiencia
    \item Al escalar a un millón de tokens o más, la ventaja en eficiencia se vuelve todavía más notable
    \item Es la \textbf{primera arquitectura que supera a full attention en todas las dimensiones} (contexto corto, entrada larga, salida larga)
\end{itemize}

La arquitectura combina linear attention con full attention en una proporción de 3:1, con lo que logra un equilibrio entre la capacidad de contexto largo y la eficiencia.

\subsection{Resumen de la sección}

Kimi Linear, mediante el factor de decaimiento de grano fino y una fórmula paralela exactamente equivalente, logra un desempeño mejor y una eficiencia mayor que full attention, con una ventaja especialmente notable en escenarios de contexto extremadamente largo.

\section{Agent Swarms: colaboración paralela multiagente}

\subsection{Paradigma de Agent Swarm}

El modo de un solo agente enfrenta un límite en la complejidad de las tareas. Kimi propuso el paradigma \textbf{Agent Swarm}:

\begin{itemize}
    \item Un \textbf{agente principal (orchestrator)} se encarga de la organización de la tarea
    \item El agente principal puede generar un conjunto de \textbf{agentes hijos} y asignarles subtareas
    \item Los agentes hijos se ejecutan en paralelo y el agente principal recopila los resultados
    \item Todo el proceso puede realizarse de forma iterable
\end{itemize}

\begin{knowledgebox}{Analogía con una organización humana}
Agent Swarm es similar a la estructura organizacional de una empresa: el CEO se encarga de descomponer y asignar tareas (orchestrator), mientras que los investigadores de IA, los desarrolladores web, los investigadores de física, etc., cumplen cada uno su función (sub-agents), y finalmente roles como el fact checker agregan los resultados.
\end{knowledgebox}

Los experimentos muestran que Agent Swarm puede reducir de manera significativa el tiempo de ejecución de las tareas, con un efecto especialmente notable en tareas de alta complejidad. Al escalar a 100 o incluso 1,000 agentes hijos, es posible completar tareas complejas en un tiempo aceptable.

\subsection{Entrenamiento con RL de Agent Swarm}

\begin{importantbox}{Tres funciones de recompensa}
Además de la recompensa de resultado (outcome reward) estándar, el entrenamiento con RL de Agent Swarm introduce dos señales de recompensa adicionales:
\begin{enumerate}
    \item \textbf{Instantiation Reward}: incentiva la instanciación paralela de agentes hijos y evita el «colapso serial», es decir, que el modelo degenere hacia la ejecución de un solo agente
    \item \textbf{Finish Reward}: asegura una tasa alta de finalización de las subtareas y evita que el modelo explote la primera recompensa, generando una gran cantidad de agentes hijos sin completar la tarea
    \item \textbf{Outcome Reward}: la recompensa estándar por grado de finalización de la tarea
\end{enumerate}
El peso de las dos primeras recompensas disminuye gradualmente durante el entrenamiento (decay strategy): al inicio se fomenta la exploración y, más adelante, se privilegian los resultados.
\end{importantbox}

\begin{warningbox}{Trampas en el entrenamiento de Agent Swarm}
Si solo se usa la outcome reward, el modelo tiende a degenerar hacia el modo de un solo agente (colapso serial). Si solo se agrega la instantiation reward, el modelo puede hacer hack de esta señal, generando una gran cantidad de subtareas falsas sin completarlas nunca. Es necesario usar las tres señales de recompensa en conjunto para lograr un entrenamiento estable.
\end{warningbox}

\subsection{Resumen de la sección}

Agent Swarm abre una nueva dimensión de scaling mediante la colaboración paralela multiagente. Las tres funciones de recompensa, cuidadosamente diseñadas y combinadas con una estrategia de decaimiento, permiten que el modelo aprenda de manera efectiva la orquestación paralela y la capacidad de descomposición de tareas.
\section{Kimi K2.5: fusión de tres dimensiones}

\subsection{Entrenamiento y arquitectura}

K2.5 fusiona los avances de las tres dimensiones anteriores en un solo modelo:
\begin{itemize}
    \item \textbf{Optimizador Muon + QK-Clip}: mejora la token efficiency
    \item \textbf{Kimi Delta Attention + arquitectura Linear}: refuerza la capacidad de long context
    \item \textbf{Agent Swarms}: abre una nueva dimensión de parallel scaling
\end{itemize}

El modelo se entrenó en GPU NVIDIA H800, con cada nodo con 2 TB de RAM y las GPU interconectadas mediante NVLink. El modelo base de K2.5 se entrenó con más de 30 billones de tokens (base model 15T + 15T adicionales de K2.5), y el proceso de entrenamiento fue extremadamente estable: no hubo ningún loss spike.

\begin{importantbox}{Early Fusion: entrenamiento conjunto nativo de visión y texto}
K2.5 es el primer modelo abierto con \textbf{capacidad nativa conjunta de visión y texto}. A diferencia de los esquemas de late fusion, que agregan capacidad visual sobre un modelo de texto ya existente, K2.5 fusiona los tokens de visión y de texto desde el día 0 del entrenamiento (early fusion); los experimentos iniciales muestran que este esquema supera al de late fusion.
\end{importantbox}

\subsection{Efecto de refuerzo entre modalidades}

Durante el entrenamiento de K2.5 se descubrió un fenómeno interesante: las dos modalidades pueden reforzarse mutuamente.

\begin{importantbox}{Vision mejora Text, Text mejora Vision}
\begin{itemize}
    \item \textbf{Vision $\rightarrow$ Text}: usar RL únicamente con tareas visuales (sin ninguna tarea de matemáticas o programación) logra, aun así, mejorar el desempeño en razonamiento textual
    \item \textbf{Text $\rightarrow$ Vision}: con una base textual sólida, no se necesitan en absoluto datos de SFT visual (Zero Vision SFT); basta con SFT textual + RL conjunto para alcanzar un desempeño visual cercano al SOTA
\end{itemize}
Este hallazgo muestra que el early fusion hace que las dos modalidades compartan de verdad el espacio de representación, lo que permite una transferencia de capacidades en ambos sentidos.
\end{importantbox}

\section{Attention Residue: avance de la arquitectura de próxima generación}

\subsection{De ResNet a Attention Residue}

Yang Zhilin (杨植麟) recordó el tutorial clásico de He Kaiming (何恺明) en ICML 2016: ResNet resolvió el problema del desvanecimiento del gradiente en el entrenamiento de redes profundas, lo que permitió entrenar redes de cualquier profundidad. Citó la opinión de Ilya Sutskever de dos años antes: \textbf{Residual connection es, en esencia, una LSTM rotada 90 grados}.

\begin{knowledgebox}{La «atención» en la dimensión de profundidad}
Si en la dimensión temporal, pasar de LSTM a Attention fue una actualización exitosa, ¿se puede hacer lo mismo en la dimensión de profundidad? Residual connection (similar a la compuerta aditiva de la LSTM) puede sustituirse por Attention: no solo toma la salida de la capa anterior, sino que agrega mediante atención \textbf{las salidas de todas las capas anteriores}. Esto es precisamente Attention Residue.
\end{knowledgebox}

\subsection{Block Attention Residue}

Para reducir el costo de comunicación y de memoria de la GPU, el equipo diseñó además \textbf{Block Attention Residue}:
\begin{itemize}
    \item Dividir todas las capas en múltiples blocks (por ejemplo, un block cada 16 capas)
    \item Dentro de cada block se usa una residual connection estándar
    \item Entre blocks se usa attention residue
\end{itemize}

\subsection{Resultados experimentales}

\begin{itemize}
    \item En el scaling law se logró una mejora del 24\% en token efficiency (50T token $\rightarrow$ equivalente a 62T token)
    \item Validation loss se mantuvo consistentemente por debajo de la curva original
    \item La mejora fue más notable en benchmarks intensivos en razonamiento como GPQA, Math y HumanEval
\end{itemize}

\subsection{Resumen de la sección}

Attention Residue trasladó a la dimensión de profundidad el éxito del paso de LSTM a Attention en la dimensión temporal, y con una mejora del 24\% en token efficiency y avances notables en tareas de razonamiento, demostró que todavía hay un enorme margen para la innovación arquitectónica.

\section{Síntesis y ampliación}

\subsection{El cambio de paradigma en la investigación}

Yang Zhilin (杨植麟) señala que el paradigma de investigación en IA actual es radicalmente distinto al de hace 10 años. Antes se privilegiaba «publicar una nueva idea», pero sin el respaldo de experimentos rigurosos era difícil llegar a conclusiones confiables. Ahora, con la scaling ladder y un rico conjunto de benchmarks, es posible validar las ideas en distintas escalas, lo que permite hacer mejoras seguras y sólidas sobre técnicas «antiguas» (como los optimizadores, el mecanismo de atención y las conexiones residuales).

\subsection{El efecto multiplicativo de las tres dimensiones}

\begin{importantbox}{De la suma a la multiplicación}
Las tres dimensiones de scaling no se suman de manera simple, sino que mantienen una relación multiplicativa de refuerzo mutuo:
\begin{itemize}
    \item Adam (2014) $\rightarrow$ Muon-Clip: mejor token efficiency
    \item Full Attention (2017) $\rightarrow$ Kimi Linear: mejor contexto largo
    \item Residual Connection $\rightarrow$ Attention Residue: scaling de profundidad más eficiente
\end{itemize}
Al multiplicar estas mejoras, se obtiene una ganancia global muy superior a la suma de las mejoras individuales.
\end{importantbox}

\subsection{El futuro de la comunidad de código abierto}

Yang Zhilin se muestra optimista sobre el futuro de la comunidad de código abierto y considera que seguirán surgiendo mejoras revolucionarias en arquitectura y optimización. Agent Swarms no es el punto final: seguirán apareciendo nuevas dimensiones de scaling, lo que impulsará una competencia continua entre los modelos de código abierto y los de código cerrado.

\subsection{Lecturas adicionales}

\begin{itemize}
    \item Kaplan et al., «Scaling Laws for Neural Language Models» (2020) --- scaling law clásica
    \item Artículo de Muon Optimizer --- trabajo del equipo de Kimi sobre un optimizador de segundo orden
    \item Reporte técnico de Kimi Linear --- derivación detallada de las fórmulas de Delta Attention
    \item Reporte técnico de Attention Residue --- experimentos completos sobre la atención en la dimensión de profundidad
    \item He et al., «Deep Residual Learning for Image Recognition» (2015) --- artículo original de ResNet
\end{itemize}

%% --- Fin del contenido --- %%

\end{document}
